\section{Linking, Gauss words and Brunnian links}
\label{sec:linking-surface}

For arithmetic statements involving planarity, take the supplied rational
polygonal realization of the diagram. Crossings are transverse intersections
in the interiors of segments. No vertex meets another strand and no three
strands meet. Label the components and direct each cyclic edge list.
Write $\epsilon(p)$ for the crossing sign, choosing the positive two-strand
braid generator to have sign $+1$ when both strands follow the braid.
Equivalently this is a fixed sign times
$\operatorname{sgn}\det(u_{\rm over},u_{\rm under})$ for the two
directed tangents. The fixed sign depends on the orientation convention
for the drawing plane and cancels in the argument below.

\begin{lemma}[Signed planar intersection]
\label{lem:rational-mixed}
For two closed oriented rational polygonal walks $P,Q$ in mutual general
position, the sum of $\operatorname{sgn}\det(u_P,u_Q)$ at their
intersections, counting pairs of passages, is zero. In particular their
number of intersections is even.
\end{lemma}
\begin{proof}
Choose a point $a$ avoiding the supporting lines of the edges of $P,Q$,
the lines from vertices of $P$ through vertices of $Q$, and the lines
through vertices of $P$ parallel to edges of $Q$. From $N$ forbidden lines
one obtains such an $a$ by testing $(j,j^2)$ for $0\leq j\leq2N$.
A line excludes at most two tests. For each successive pair $p_i,p_{i+1}$
use the directed boundary of the triangle $(a,p_i,p_{i+1})$.
For a counterclockwise triangle the summand is $+1$ when $Q$ enters
and $-1$ when it leaves, since its interior lies to the left of each edge.
The sum telescopes to zero along closed $Q$. One may split its edges at
the finitely many intersections and use the inside bits of the consecutive
open pieces. A clockwise triangle multiplies this identity by $-1$.
Adding the identities cancels each spoke in opposite directions and
leaves precisely the intersection sum for $P$. Equal positive and negative
counts give the parity assertion.
\end{proof}

\begin{theorem}[Three crossing formulas for linking]
For distinct components $i,j$, the signed sum $O_{ij}$ at crossings
where $i$ passes over $j$ equals $O_{ji}$. Their common value
$\ell_{ij}$ satisfies
\[
  \ell_{ij}=\ell_{ji}
  =\frac12\sum_{p\text{ between }i,j}\epsilon(p).
\]
It is invariant along every Reidemeister movie preserving the component
labels and orientations. Mirroring negates it. Reversing exactly one
of the two components negates it, and reversing both preserves it.
\end{theorem}
\begin{proof}
In the intersection sum ordered as $(i,j)$ the over-$i$ terms agree with
their crossing signs up to the fixed sign, and the over-$j$ terms have the
opposite comparison. Lemma~\ref{lem:rational-mixed} gives
$O_{ij}-O_{ji}=0$. A curl involves only one component. In a type II
bigon the same strand is upper at the two crossings and the directed
intersection signs are opposite. Indeed, following one side of the bigon
across the other strand enters and exits on opposite sides. Reversing
either orientation flips both local comparisons. In type III, each
pair of participating strands crosses once on either side, with the same
relative height, directions and crossing sign. These verifications also
apply when two of the participating strands belong to the same component.
An over/under switch reverses the determinant order. Reversing one tangent
changes its sign. The transformation laws follow term by term.
\end{proof}

Let $w$ be the sum of all crossing signs and $w_i$ the sum over crossings
of $i$ with itself. Thus
\[
  w=\sum_i w_i+2\sum_{i<j}\ell_{ij}.
\]
The same local proof leaves $w$ fixed in types II and III. A curl changes
it by its sign. Mirroring negates $w$ and reversing all components
preserves it. For any diagram, the number of crossings between $i,j$
is at least $2|\ell_{ij}|$. A change of such a crossing changes
$\ell_{ij}$ by exactly one in absolute value. A crossing change involving
another unordered pair, or a self-crossing, leaves it fixed. Consequently
every movie using ordinary moves and crossing changes to separate a
partition $S\mid T$ of the components into disjoint diagram boxes
uses at least $\sum_{i\in S,j\in T}|\ell_{ij}|$ crossing changes.
This follows from the triangle inequality at each step, since all these
entries vanish in the separated boxes.

\begin{corollary}[Even two-strand torus links]
For the closure of $\sigma_1^{2r}$, $r\in\mathbb Z$, the components with
their braid orientations have linking number $r$. The least crossing count
in its diagram-movie class is $2|r|$, and the least number of crossing
changes needed to obtain the two-component unlink is $|r|$.
In particular the Hopf link has crossing number two.
\end{corollary}
\begin{proof}
The strand permutation is the identity. Every one of the $2|r|$
crossings is between the two components and they have the same sign
when $r\ne0$. This realizes the crossing lower bound. Switch one letter
in each consecutive like-signed pair; each pair then cancels by type II.
The resulting empty two-braid closes to the unlink. The linking bound
gives the reverse inequality, including $r=0$.
\end{proof}

\begin{theorem}[The Dowker even--odd rule]
In a one-component plane diagram with $c$ crossings, label the $2c$
successive visits $1,\ldots,2c$ starting at any regular basepoint.
The two labels at every crossing have opposite parity. Hence in order
of the odd labels the partner labels are a permutation of $2,4,\ldots,2c$.
\end{theorem}
\begin{proof}
Cut the parametrizing circle at its two visits to a particular crossing.
Each of the two resulting arcs closes locally: its two ends lie on
different strands and hence at adjacent corners of a crossing box.
The two reconnections can be drawn separately along opposite sides of
that box. Between the two visits each self-crossing of the chosen arc
occurs twice and each crossing with the other arc occurs once. The
latter number is even by Lemma~\ref{lem:rational-mixed} applied to
the reconnected polygons. If the visits have labels $a<b$, their
difference is one more than this even total $b-a-1$. Each even visit
therefore partners one odd visit. Ordering by the latter yields the
claimed permutation; recording whether the even visit is under or over
provides the extra sign of a signed Dowker code.
\end{proof}

\begin{proposition}[The Borromean deletion test and failure of pairwise detection]
The standard Borromean closed braid $B=(\sigma_1\sigma_2^{-1})^3$ has
three components. Each two-component sublink is the unlink and all three
pairwise linking numbers vanish. The full link is distinct from the
three-component unlink under Reidemeister movies.
\end{proposition}
\begin{proof}
Label the top strands $1,2,3$. The unordered pairs exchanged successively
are $12,13,23,12,13,23$, with signs $+,-,+,-,+,-$. The final ordering
is $123$, so each strand closes to its own component. After deleting any
one component, the other two cross once positively and once negatively.
Their two-strand word reduces by one cancellation and their closure is an
unlink. To distinguish all three together, give the top strands Fox colors
$(a,b,c)$ modulo $q\ge2$. A positive letter sends the affected pair
$(x,y)$ to $(2x-y,x)$ and a negative letter to $(y,2y-x)$. A closure
coloring is exactly a fixed vector. Substituting the six letters gives
\[
 (M-I)\begin{pmatrix}a\\b\\c\end{pmatrix}
 =4\begin{pmatrix}2&-1&-1\\-1&0&1\\-3&1&2\end{pmatrix}
   \begin{pmatrix}a\\b\\c\end{pmatrix}.
\]
The second row and then the first give $4(c-a)=4(b-a)=0\pmod q$;
these conditions imply all three rows. Each difference has
$\gcd(4,q)$ choices, since $4x=0\pmod q$ means
$q/\gcd(4,q)$ divides $x$. Thus the color count is
$q\gcd(4,q)^2$. At $q=3$ this gives $3$, whereas the unlink has
$3^3=27$. The Reidemeister bijections for Fox colorings prove the
distinction. In particular, even the complete list of pairwise linking
numbers cannot detect all linking.
\end{proof}


\begin{corollary}[Six crossings are necessary for the Borromean rings]
The least crossing count in the Borromean diagram-movie class is six.
\end{corollary}
\begin{proof}
Use the bracket and span bound of Section~\ref{sec:bracket}. In the
six-letter braid above, the writhe is zero. Here is its complete state
count grouped by the number $b$ of $B$-pairings and the number $h$ of
circles. The rows are $b=0,\ldots,6$ and the columns are $h=1,2,3,4$:
\[
  \begin{pmatrix}
  0&0&0&1\\0&0&6&0\\0&15&0&0\\16&0&4&0\\
  0&15&0&0\\0&0&6&0\\0&0&0&1
  \end{pmatrix}.
\]
To reproduce each entry, trace three columns through seven horizontal
levels. At letter $j$ join each unaffected column to its successor.
Pair the two affected columns vertically for an $A$ at a positive
letter or a $B$ at a negative letter. Otherwise join the two upper
entries and join the two lower entries. Identify bottom and top and
count equivalence classes on the $21$ entries. Summing
$A^{6-2b}(-A^2-A^{-2})^{h-1}$ with the displayed multiplicities, and
putting $X=A^4$, gives
\[
 -X^3+3X^2-2X+4-2X^{-1}+3X^{-2}-X^{-3}.
\]
Its value at $X=-1$ is $16$. A disconnected shadow forces the bracket
to have a factor $-A^{-2}(A^4+1)$, by splitting its nonempty pieces
in the state sum. Hence every equivalent diagram has connected shadow.
The $A$-span is $24$, so the bound $\operatorname{span}_A\le4c$ gives
$c\ge6$. The braid attains equality.
\end{proof}

\section{A rational Seifert surface and its genus certificate}
\label{sec:seifert}
At a crossing there are two incoming and two outgoing corners. Each
through strand joins opposite corners, so the two incoming corners
are neighbors in cyclic order. Replace the crossing by the unique
noncrossing pairing joining incoming to outgoing. The result consists of
$s$ disjoint oriented polygonal circles, the \emph{Seifert circles}.
Isolated crossingless circles are included. The following elementary
capping lemma makes the geometric content of the disk-and-band construction
explicit.

\begin{lemma}[Planar polygon caps]
\label{lem:slab-cap}
There is a primitive recursive construction of a rational triangulated
disk bounded by a supplied simple rational polygon $P$. For finitely
many disjoint such polygons their planar disks have disjoint interiors
or are nested, and the strict containments are decidable.
\end{lemma}
\begin{proof}
Shear by $(x,y)\mapsto(x+hy,y)$ so that the vertices have distinct
first coordinates. Each vertex pair forbids at most one integer $h$.
Split the plane at these coordinates into vertical slabs. In an open slab,
sort the intersecting edges by height at the midpoint abscissa. This order
is constant within the slab, since edges do not cross. The gaps between
edges have alternating labels, starting with $0$ above the polygon.
Across an open piece of a slab wall the labels agree: at a vertex the
number of edges above changes by zero or two. Thus crossing any polygon
edge changes the label, and walking through a wall away from the polygon
preserves it. The label-one gaps are bounded triangles or trapezoids.

Here are the two connectivity facts needed about this subdivision. On a
small punctured neighborhood of a vertex, its two incident segments cut
the neighborhood into two sectors. The labels differ across the segments,
so the label-one sides of its consecutive edges meet through the same
sector. Walking the boundary in order therefore connects every edge's
label-one side to every other. Every label-one gap has a polygon edge
as its ceiling or floor, giving connectedness of the label-one cells by
wall adjacency. The same argument applies to the label-zero sides.
Every bounded label-zero gap touches such a side, and the outside gaps
join above, below or around an extreme vertex. Consequently every
label-zero point away from $P$ can be routed to the unbounded gap.
These statements concern finite cell adjacency. Near a vertex one can
take a square of radius less than its distances to nonincident segments
and other vertices.

Use the closures of the label-one gaps. Subdivide both copies of a
common wall at all polygon points on it. Triangulate a convex
trapezoid by a diagonal, with a single triangle when one end collapses.
To accommodate collinear marks on a side, split its incident triangle
from the opposite vertex, successively. All vertices of the resulting
triangles lie on $P$: even a wall endpoint in the middle of a polygon
edge lies on that edge. The dual graph, with one vertex per triangle
and one edge across each interior side, is connected by the preceding
cell argument and the triangulation within each gap.

Every interior side $e$ is a bridge in this graph. Join its endpoints
by either boundary subarc of $P$ to form a closed polygon with $e$.
If a dual path could bypass its dual edge, the bypass together with
that dual edge would give a closed walk through triangle interiors
crossing $e$ exactly once and disjoint from $P$. To draw the walk use
triangle interior points and side midpoints. At a crossed side join
nearby interior points directly so that its intersection is transverse
and interior to a segment. This contradicts the parity assertion of
Lemma~\ref{lem:rational-mixed}. The dual graph is thus a tree.
The triangulation can be assembled by attaching a triangle on a single
free side at each step along the tree. No extra identification at an
isolated corner is needed: around each boundary vertex the incident
triangles form one consecutive fan in its interior sector. Hence all
copies of that corner are already identified through shared sides.
A triangle attached on one entire boundary side to a disk gives a disk
(replace that side of the boundary by the other two), proving the claim.
The tree order itself is a disk certificate.

For disjoint polygon boundaries $P,Q$, the bit for $P$ is constant
along $Q$. If $Q$ has bit one, the exterior of $P$, together with its
boundary, has routes to infinity avoiding $Q$. These points are exterior
to $Q$; its disk therefore lies inside the disk of $P$. If each boundary
is exterior to the other, the bit for $Q$ is constant on the connected
interior of $P$ and is zero just inside its boundary, hence the disk
interiors are disjoint. Both bits cannot be one, by strict containment.
A wall or ray meeting a vertex is handled by its neighboring gap, or by
subdivision and an interior test point. Inverting the shear restores the
original rational coordinates.
\end{proof}

\begin{theorem}[Seifert construction with a boundary witness]
\label{thm:seifert-pl}
From a nonempty oriented rational plane link diagram with $c$ crossings,
$b$ components and $s$ Seifert circles, one constructs by primitive
recursion a rational embedded orientable triangulated surface whose
boundary projects to that same decorated diagram, up to subdivision.
It has $k$ connected components, where $k$ is the number of components
of the incidence graph of Seifert circles and crossing bands. Its Euler
characteristic and total handle genus $G$ satisfy
\[
  \chi=s-c,\qquad 2G=2k-b-s+c.
\]
In particular a knot diagram gives a connected surface of genus
$g_D=(1+c-s)/2$, an explicit upper bound for the spanning genus of the knot.
\end{theorem}

\begin{proof}
Fit disjoint rational parallelograms at the crossings, with corners
on the four strand ends. They can be chosen as the convex hull of
$p\pm t u,p\pm t v$ at a crossing $p$ of strands with directions
$u,v$. Take pairwise crossing distances, distances from crossings to
nonincident closed segments and to the endpoints of their incident
segments. Their positive squared distances have a rational minimum $m$
(append $1$ before taking minima). For a point and a segment the squared
distance is computed by projecting to its supporting line and clamping
the rational parameter to $[0,1]$. With $H$ the maximum of $1$ and
the $\ell^1$ norms of the directions, one may take
$t=\min(1,m)/(100H)$. Thus no unrelated pieces enter these
parallelograms. Their opposite sides used by the oriented smoothing
serve as the two attachment intervals. Outside them retain the input
segments. This gives the circles to which Lemma~\ref{lem:slab-cap}
applies.

Let $d(C)$ be the number of other disks strictly containing circle $C$,
and put its disk at height $-s+d(C)<0$. Connect its rim vertically
to the copy of $C$ at height zero using a triangulated wall of
rectangles. A containing disk lies below every contained one. The walls
of different circles have disjoint projections. A wall above an inner
disk never reaches the level of its containing cap, and disjoint
nonnested disks have disjoint projections. This proves pairwise
disjointness of the capped walls. Each is a disk with a collar and is
oriented to induce the specified direction on its upper boundary.

Use affine horizontal coordinates in each parallelogram so that the
attachment intervals are $x=-1$ and $x=1$ in the square $[-1,1]^2$.
The smoothing directs both intervals the same way in $y$. Here is a
specific half-twisted PL band. One boundary side visits
\[
 (-1,-1,0),\quad (0,0,1),\quad (1,1,0),
\]
and the other visits $(-1,1,0),(0,0,3),(1,-1,0)$.
Connect the two points in each column $x=-1,0,1$ and split each of the
two resulting four-sided patches along a diagonal into two triangles.
Within either half, one triangle projects onto a triangle of positive
area and the other onto a boundary segment of that triangle; their
intersection is precisely their common edge. The two halves meet only
on the middle vertical edge. The band meets the plane $z=0$ exactly
at its attaching ends. Its sides project to the original diagonals,
with the specified overpass obtained by interchanging the heights
$1,3$ when necessary. Because the order of the two side labels reverses
between the ends, the induced boundary directions of the rectangle
along both attachments point the same way in the drawing. Orient it
oppositely to the two disk boundary intervals. This consistently
orients the union. The patches occupy distinct boxes, and the rest of
the surface lies at $z\leq0$. The resulting list of triangles is an
embedded surface. Boundary subdivisions of the attachment intervals
are matched by splitting the corresponding band triangles from their
opposite corners. Its boundary follows the original crossing assignment
at every box and the original edges outside.

The disks and bands also give a purely finite computation of genus.
Initially there are $s$ disks with $s$ boundary circuits. Attaching a
rectangle on two disjoint intervals reduces the Euler characteristic
by one: after a common subdivision it contributes $1-2$ by
inclusion--exclusion for vertex, edge and face counts. The graph
components track the surface components. For a prefix of the bands let
$k',b',\chi'$ be the component, boundary and Euler counts. An oriented
band on two intervals of the same boundary circle splits that boundary
in two (trace the two complementary arcs after removing the intervals).
On two distinct boundary circles it merges them into one. In the latter
case the surface components also merge exactly when the intervals
belonged to distinct components. Starting with $G'=0$, increment $G'$
only for a band joining distinct boundary circles of the same component.
These three updates preserve
\[
                      2G'=2k'-b'-\chi'.
\]
They are also the handle description of orientable genus: a band on
one boundary splits a collar into a pair of pants; a band joining two
boundaries in one connected surface makes a handle, merging those
boundaries. Equivalently the displayed Euler relation computes total
genus for an orientable surface with boundary. At completion the
boundary circuits trace the diagram and number $b$. Every component of
a disk-and-band surface has boundary, by the same updates, so $b=1$
forces $k=1$. This proves all the formulas.
\end{proof}

\begin{corollary}[Bounds and operations]
A knot diagram with $c$ crossings bounds a surface of genus at most
$\lfloor c/2\rfloor$. For a connected sum made at regular arcs of two
separate oriented knot diagrams, the surface genus of the resulting
diagram is $g_{D\#E}=g_D+g_E$. Mirroring preserves $g_D$. For $p,q\ge1$,
the usual closure of $(\sigma_1\cdots\sigma_{p-1})^q$ has
$\gcd(p,q)$ components and a connected spanning surface of genus
\[
   \frac{2-\gcd(p,q)-p+(p-1)q}{2}.
\]
For coprime $p,q$ this gives the knot-genus upper bound
$(p-1)(q-1)/2$.
\end{corollary}
\begin{proof}
Use $s\ge1$ in the knot formula. In a connected sum, smoothing all
crossings on either side leaves exactly the two marked Seifert circles
spliced into one, so the crossing counts add and the circle counts add
with a subtraction of one. A mirror changes heights without changing
the oriented smoothing. For the displayed braid, with all strands
oriented along the braid, smoothing leaves all $p$ columns vertical and
each closes to its own circle. There are $(p-1)q$ crossing bands.
The incidence graph contains the path through adjacent columns and is
connected (also for $p=1$). The permutation of one block rotates the
$p$ labels by one, up to reverse convention. Its $q$th power has
$\gcd(p,q)$ orbits, since the subgroup of residues generated by $q$
consists of the multiples of $\gcd(p,q)$ modulo $p$, by the Euclidean
algorithm. Substitute these counts. Each upper bound is witnessed by the
constructed surface. Minimality among spanning surfaces would require a
lower bound and is not used in the construction.
\end{proof}

\paragraph{Arithmetic and Lean implementation.}
The planar-intersection loop stores the two signed counts and the current
triangle-inside bit. Its induction invariant is that the signed prefix
equals the final minus initial bit, multiplied by the triangle orientation.
For a polygon with $n$ vertices the shear has index at most $n(n-1)/2$.
There are at most $n+1$ slabs and at most $n+1$ gaps per slab.
Each slab wall has at most $n$ intersections with polygon edges plus
its vertex. Thus $10(n+1)^3$ suffices for the triangles of a single cap
and their intervening wall subdivisions. Connected components of a graph
with $v$ vertices may be computed by $v$ iterations of adding all
neighbors to a bit vector. The invariant says that its bits are exactly
the endpoints of walks of length at most the stage; walk certificates can
be bounded by $(v+1)^{v+1}$ for these stages. A cycle or bypass can be
shortened to a simple one by deleting repeated-vertex portions, with at
most $v+1$ successive retained vertices. Every instance of connectivity
in the cap proof uses these finite adjacency tests. The triangle-tree
gluing order gives local side and vertex fan checks. The depth list is
an $s$-by-$s$ table of ray parities, and the $c$ band additions carry
the equality $2G'=2k'-b'-\chi'$ as a signed-integer equality at every
prefix. There are at most $b+c$ Seifert circles, since each smoothing
changes the number of circuits by one in absolute value. All scalar
geometry uses rational determinants, additions, products and divisions
with nonzero denominators checked before use.

For Lean, implement a closed walk as a cyclic vertex array, a triangle
as a triple of rational pairs, and a surface as a list of rational
spatial triangles with oriented edge incidences. First prove triangle
entry/exit telescoping, then spoke cancellation and the signed crossing
identities. Encode a slab cell by its two adjacent abscissas and its
ordered floor/ceiling edges. The label is its vertical index modulo two;
prove wall agreement by inspecting the two edges at the wall vertex.
Prove the tree claim by converting a simple dual cycle into the rational
polygon used in Lemma~\ref{lem:slab-cap}. Independently encode each
disk rim by a cyclic word in its marked intervals. On band addition,
split a word at two marks or splice two words, removing the attaching
intervals and inserting the long band sides. The local cases above
prove the genus recurrence by induction on the word of bands.
Relate this combinatorial boundary to the triangle boundary by checking
the single two-patch band and the capped-wall template, then concatenate
the incidences. For object-theory induction, equality of arrays,
validity of all incidences and correctness of all bounded labels are
Boolean folds. This leaves only quantifier-free induction with the
input code as a free parameter.
